\documentclass[12pt,a4paper]{article}
\usepackage[utf8]{inputenc}
\usepackage[T1]{fontenc}
\usepackage[ngerman]{babel}
\usepackage[margin=2cm,top=3.5cm,bottom=2.5cm]{geometry}
\usepackage{enumitem}
\usepackage{setspace}
\usepackage{parskip}
\usepackage{fancyhdr}
\usepackage{ragged2e}
\usepackage{tikz}


\title{}
\author{}
\date{}

\setlength{\parindent}{0pt}
\setlength{\parskip}{6pt}
\renewcommand{\baselinestretch}{1.15}

\newcommand\Committee{SEDE}
\newcommand\Fraktion{LEER}

\newcommand\Thema{Armee}
\newcommand\city{Ulm}
\newcommand\datum{11.05.2026}
\newcommand\timeVorb{07:30-09:00}
\newcommand\timeEinf{09:00-09:30}
\newcommand\timeFrakOne{09:30-11:00}
\newcommand\timePauseOne{11:00-11:15}
\newcommand\timeAuss{11:15-12:30}
\newcommand\timeMittag{12:30-13:00}
\newcommand\timeFrakTwo{13:00-13:30}
\newcommand\timeFinVerh{13:30-14:00}
\newcommand\timePlenar{14:00-15:15}
\newcommand\timeDebr{15:15-15:30}
\newcommand\politiker{Andrea Wechsler}
\newcommand\politikerOffice{Mitglied des Europäischen Parlaments}
\newcommand\stadtvertreter{Nalan Schmidt}
\newcommand\stadtvertreterOffice{Anlauf- und Koordinierungsstelle Bürgerdialog (Open Government) der Stadt Ulm}
\newcommand\localSupport{das Europe Direct Ulm}
\newcommand\sponsor{die Vertretung der Europäischen Kommission in München}
\newcommand\jefvorsitz{Farras Fathi}
\newcommand\gendervorsitz{Landesvorsitzender}
\newcommand\evpLeader{Tobias}
\newcommand\evpRoom{Kleiner Sitzungssaal}
\newcommand\sdLeader{Antonia}
\newcommand\sdRoom{Veranstaltungsraum Europe Direct}
\newcommand\reLeader{Alica}
\newcommand\reRoom{Besprechungszimmer Donaubüro I}
\newcommand\fifthLeader{Farras}
\newcommand\fifthRoom{Besprechungszimmer Öffentlichkeitsarbeit}
\newcommand\pfeLeader{Hannah}
\newcommand\pfeRoom{Besprechungszimmer BM III}
\newcommand\numSuS{27}
\newcommand\location{im Ulmer Rathaus}
\newcommand\fifthGroup{Linke}
\newcommand\anzahlcomm{3}
\newcommand\totcoms{26}
\newcommand\minmems{25}
\newcommand\maxmems{90}

\newcommand{\LogoBasePath}{../Logos}
% Shared logo helpers for SimEP LaTeX documents.

\providecommand{\LogoBasePath}{Logos}
\providecommand{\FraktionsLogoBasePath}{Folien/Bilder}

\providecommand{\constrainedlogo}[4][2cm]{%
  \raisebox{#3}{\includegraphics[width=#2,height=#1,keepaspectratio]{#4}}%
}

\providecommand{\jeflogo}{\constrainedlogo[2.5cm]{5cm}{0pt}{\LogoBasePath/jef.png}}
\providecommand{\simeplogo}{\constrainedlogo[2.5cm]{5cm}{0pt}{\LogoBasePath/simep.png}}
\providecommand{\financelogo}{\constrainedlogo[1.5cm]{4cm}{0pt}{\LogoBasePath/finance\_\city.png}}
\providecommand{\citylogo}{\constrainedlogo[1.5cm]{4cm}{0pt}{\LogoBasePath/\city.png}}
\providecommand{\fraktionslogo}{\constrainedlogo[2.5cm]{5cm}{0pt}{\FraktionsLogoBasePath/\Fraktion.png}}

\providecommand{\PlaceLogosOne}{%
\begin{tikzpicture}[remember picture, overlay]
        % Top left corner
        \node at ([xshift=10mm, yshift=-8mm] current page.north west) [anchor=north west, inner sep=0pt] {
            {\jeflogo}
        };

        % Top right corner
        \node at ([xshift=-7.5mm, yshift=-6mm] current page.north east) [anchor=north east, inner sep=0pt] {
            {\simeplogo}
        };


        % Bottom left corner
        \node[anchor=south west, inner sep=0pt, text width=4.5cm, align=left] at ([xshift=10mm, yshift=8mm] current page.south west) {
            \small Auftraggeber \& \newline Kooperationspartner:\\[3mm]
            \financelogo
        };


        % Bottom right corner
        \node[anchor=south east, inner sep=0pt, text width=4.5cm, align=right] at ([xshift=-10mm, yshift=10mm] current page.south east) {
            \small Unterst\"utzt durch:\\[3mm]
            \citylogo
        };

\end{tikzpicture}
}

\providecommand{\PlaceLogosRest}{%
\begin{tikzpicture}[remember picture, overlay]
        % Top left corner
        \node at ([xshift=10mm, yshift=-8mm] current page.north west) [anchor=north west, inner sep=0pt] {
            {\jeflogo}
        };

        % Top right corner
        \node at ([xshift=-7.5mm, yshift=-6mm] current page.north east) [anchor=north east, inner sep=0pt] {
            {\simeplogo}
        };

\end{tikzpicture}
}

\providecommand{\PlaceLogosFraktion}{%
\begin{tikzpicture}[remember picture, overlay]
        % Top left corner
        \node at ([xshift=10mm, yshift=-8mm] current page.north west) [anchor=north west, inner sep=0pt] {
            {\jeflogo}
        };

        % Top right corner
        \node at ([xshift=-7.5mm, yshift=-6mm] current page.north east) [anchor=north east, inner sep=0pt] {
            {\simeplogo}
        };

        % Upper-right caucus logo
        \node at ([xshift=-12mm, yshift=-35mm] current page.north east) [anchor=north east, inner sep=0pt] {
            {\fraktionslogo}
        };

        % Bottom left corner
        \node[anchor=south west, inner sep=0pt, text width=4.5cm, align=left] at ([xshift=10mm, yshift=8mm] current page.south west) {
            \small Auftraggeber \& \newline Kooperationspartner:\\[3mm]
            \financelogo
        };


        % Bottom right corner
        \node[anchor=south east, inner sep=0pt, text width=4.5cm, align=right] at ([xshift=-10mm, yshift=10mm] current page.south east) {
            \small Unterst\"utzt durch:\\[3mm]
            \citylogo
        };

\end{tikzpicture}
}

\begin{document}

\thispagestyle{empty}

\vspace*{2cm}

\begin{center}
\Large
Vorschlag für eine

\vspace{0.8cm}

\textbf{\large Entschließung des Europäischen Parlaments\\[0.3em]
zur Einführung einer Europäischen Armee}

\vspace{2cm}

\normalsize
Federführende Ausschüsse:

\vspace{0.5cm}

\textbf{Haushaltsausschuss (BUDG)}

\textbf{Ausschuss für Bürgerliche Freiheiten, Justiz und Inneres (LIBE)}

\textbf{Ausschuss für Sicherheit und Verteidigung (SEDE)}
\end{center}

\PlaceLogosOne


\newpage

\subsection*{DAS EUROPÄISCHE PARLAMENT,}

\begin{itemize}[leftmargin=2em, itemsep=8pt]
    \item gestützt auf die Verträge, insbesondere auf Artikel~24, Artikel~42 und Artikel~46 des
    Vertrages über die Europäische Union (EUV) und die damit einhergehende Etablierung
    einer Gemeinsamen Außen- und Sicherheitspolitik (GASP),

    \item unter Hinweis auf die gemeinsame Mitteilung von 23 Mitgliedstaaten vom
    13.~November 2017 über die Ständige Strukturierte Zusammenarbeit in der Sicherheits-
    und Verteidigungspolitik (SSZ/PESCO) nach Artikel~42 Absatz~6 EUV,

    \item unter Hinweis auf die Gemeinsamen EU-NATO-Erklärungen vom Juli 2016, vom
    10.~Juli 2018 sowie vom 10.~Januar 2023,
\end{itemize}

\subsubsection*{IN ERWÄGUNG NACHSTEHENDER GRÜNDE: (Alle Ausschüsse)}

\begin{enumerate}[label=\Alph*., leftmargin=2em, itemsep=8pt]
    \item in der Erwägung, dass die zunehmenden militärischen Konflikte in geographisch
    teils unmittelbarer Nähe der Außengrenzen der Union eine Beschleunigung der in
    Artikel~24 Absatz~1 EUV vorgesehenen schrittweisen Festlegung einer gemeinsamen
    Verteidigungspolitik erfordern, die nach dem Vertragstext zu einer gemeinsamen
    Verteidigung führen kann;

    \item in der Erwägung, dass die existierende Strukturierte Zusammenarbeit in der
    Sicherheits- und Verteidigungspolitik (SSZ-PESCO) dringend weiter ausgebaut werden
    muss, um im Verteidigungsfall eine effektive gemeinsame Verteidigung des
    Unionsgebiets zu ermöglichen;

    \item in der Erwägung, dass die Schwierigkeiten der einzelnen Mitgliedstaaten bei der
    Beschaffung militärischen Materials, wie sie derzeit im Ukraine-Konflikt deutlich
    werden, eine Verstärkung der gemeinsamen Beschaffungspolitik erfordern;

    \item in der Erwägung, dass nur eine gemeinsame Europäische Armee mit einheitlicher
    Kommandostruktur, einheitlichem Dienstrecht und einheitlichem Zugang für alle
    Unionsbürgerinnen und -bürger die gemeinsame Verteidigung wirksam bewältigen kann,
    auch wenn derzeit noch Zuständigkeitsprobleme bestehen;

    \item in der Erwägung, dass gemeinsame Verteidigungseinrichtungen die Solidarität der
    Unionsbürgerinnen und -bürger und den Austausch untereinander zu fördern geeignet
    sind;
\end{enumerate}

\newpage

\noindent\ldots fordert im Rahmen der Gemeinsamen Außen- und Sicherheitspolitik (GASP) die Einführung einer gemeinsamen Europäischen Armee, welche folgendermaßen ausgestaltet werden soll:

\subsection*{Gegenstand und Ziele (SEDE)}

\begin{enumerate}[leftmargin=2em, itemsep=10pt]
    \item Jeder EU-Mitgliedstaat soll sich verpflichten, einen Teil seiner nationalen Streitkräfte und Ausrüstung in eine gemeinsame Europäische Armee zu übertragen, mit dem langfristigen Ziel der kompletten Auflösung aller nationaler Streitkräfte;

    \item Jeder Mitgliedstaat soll bis
    \begin{enumerate}[label=\roman*), leftmargin=1.5em, itemsep=1pt]
        \item 2030 \quad 10~Prozent,
        \item 2040 \quad 25~Prozent,
        \item 2050 \quad 50~Prozent,
        \item 2070 \quad 100~Prozent
    \end{enumerate}
    seiner Streitkräfte und militärischen Güter in die Europäische Armee übertragen.

    \begin{enumerate}[label=\alph*), leftmargin=1.5em, itemsep=6pt]
        \item Die Befehls- und Kommandogewalt obliegt dem Hohen Vertreter der EU für Außen- und Sicherheitspolitik als Oberbefehlshaber.

        \item Die Budgethoheit über die gemeinsamen Streitkräfte obliegt dem Europäischen Parlament. Der Ausschuss für Sicherheit und Verteidigung (SEDE) soll mit der Kontrolle der Streitkräfte betraut werden.

        \item Alle Einsätze der Europäischen Armee sollen der Zustimmung des Europäischen Parlaments bedürfen. Nur bei Gefahr im Verzug soll ein Einsatz ohne vorherige Zustimmung möglich sein, in welchem Fall die Zustimmung des Parlaments nachgeholt werden muss.

        \item Die Zustimmung zu einem Einsatz soll auf maximal ein Jahr begrenzt werden und soll mit Ablaufen dieser Frist verlängert werden müssen. Ausgenommen hiervon soll die direkte Verteidigung des EU-Territoriums sein.
    \end{enumerate}

    \item \ldots regt an, dass die Europäische Armee zu folgenden Zwecken eingesetzt wird:
    \begin{enumerate}[label=\alph*), leftmargin=1.5em, itemsep=6pt]
        \item Zivil-militärische Rettungseinsätze und Katastrophenhilfe in den EU-Mitgliedstaaten sowie außerhalb der EU;

        \item Militärische Beratung und Unterstützung für die bestehenden Armeen der EU-Mitgliedstaaten;

        \item Verteidigungszwecke gegen Bedrohungen der europäischen Sicherheit;

        \item Humanitäre Interventionen in Krisengebieten zur Konfliktbewältigung und Stabilisierung
    \end{enumerate}
\end{enumerate}

\subsection*{Finanzierung (BUDG)}

\noindent Das Europäische Parlament \ldots

\begin{enumerate}[leftmargin=2em, start=4, itemsep=10pt]
    \item \ldots schlägt vor, die Finanzierung der Europäischen Armee wie folgt zu gestalten:
    \begin{enumerate}[label=\alph*), leftmargin=1.5em, itemsep=6pt]
        \item Die Europäische Kommission soll die Ausgaben für die Gemeinsame Außen- und Sicherheitspolitik im nächsten Haushaltsplan um 20\% erhöhen.

        \item Die Mitgliedstaaten sollen sich anteilig an der Finanzierung eines Fonds zum Aufbau einer Europäischen Armee beteiligen. Der jährliche Beitrag soll hierbei 0,2~\% des BIP pro Kopf der Mitgliedstaaten betragen.

        \item Während zwei Drittel des Fonds für den Aufbau der Armee genutzt werden sollen, soll ein Drittel in die militärische Forschung und den Ausbau der Europäischen Verteidigungsagentur fließen sowie der Schaffung von Arbeitsplätzen in der Europäischen Rüstungsindustrie dienen.
    \end{enumerate}

    \item \ldots schlägt vor, dass Projekte im Verteidigungsbereich von Mitgliedstaaten grundsätzlich EU-weit auszuschreiben sind. Beschränkungen, die das Prinzip der Gleichheit verletzen und Unternehmen aus dem betroffenen Mitgliedstaat bevorzugen, sind unzulässig. Über Ausnahmen entscheidet die Europäische Kommission im Einzelfall.

    \item \ldots schlägt vor, dass gemeinsame Beschaffungsprojekte im Rüstungsbereich mit bis zu 5\% des (Netto-) Auftragsvolumens gefördert werden können. Darunter fallen ebenso Rüstungsprojekte in Kooperation mit weiteren NATO-Mitgliedstaaten und der Ukraine.

    \item \ldots schlägt vor, dass die Europäische Verteidigungsagentur beauftragt wird, gemeinsame Standards für militärische Ausrüstung zu entwickeln. Diese dürfen bereits existierenden Standards der North Atlantic Treaty Organization (NATO) nicht widersprechen.
\end{enumerate}

\subsection*{Dienstrecht (LIBE)}

\noindent Das Europäische Parlament \ldots

\begin{enumerate}[leftmargin=2em, start=8, itemsep=10pt]
    \item \ldots regt an, dass die Union ein einheitliches Dienstrecht für die Berufs- und Zeitsoldatinnen und -soldaten sowie zivile Mitarbeitende der ihr zur Verfügung gestellten Truppenteile schafft, wobei langfristig eine gleiche Besoldung, Versorgung und disziplinarische Behandlung aller Militärangehörigen der Union angestrebt wird.

    \item \ldots fordert, dass neue Soldatinnen und Soldaten und zivile Mitarbeitende der durch die Mitgliedstaaten überlassenen oder direkt von der Union aufgestellten Streitkräfte direkt bei Einstellung dem zu schaffenden Dienstrecht unterstellt werden.

    \item \ldots regt an, die Laufbahn als europäische Berufs- oder Zeitsoldatin oder -soldat sowie als zivile Mitarbeitende der Streitkräfte der Union für Frauen und Männer attraktiv auszugestalten und aktiv zu bewerben.

    \item \ldots regt an, Regelungen für eine verpflichtende, einjährige Wehrpflicht bei den Streitkräften der Europäischen Union im ersten Jahr der Volljährigkeit für alle Unionsbürgerinnen und -bürger zu schaffen. Somit wird die Wehrfähigkeit der Union sichergestellt. Das Recht, den Dienst an der Waffe zu verweigern, ist dabei jedoch unbedingt anzuerkennen.
\end{enumerate}


\end{document}